%\documentclass[12pt]{article}

\documentclass[english,12pt]{article}
\usepackage[T1]{fontenc}
\usepackage{lmodern}
\usepackage[utf8]{inputenc}

\usepackage{geometry}
\usepackage{epsfig}
\usepackage{cite}
\usepackage{color}
\usepackage{amsmath}
\usepackage{amssymb}
\usepackage{multirow}
\usepackage{bigdelim}
\usepackage{microtype}
\usepackage{hyperref}
\usepackage{booktabs}
\usepackage{longtable}
\usepackage{soul}

\newcommand\sss{\mathchoice%
{\displaystyle}%
{\scriptstyle}%
{\scriptscriptstyle}%
{\scriptscriptstyle}%
}
\def\lq{\left[}   
\def\rq{\right]} 
\def\rg{\right\}} 
\def\lg{\left\{} 
\def\({\left(} 
\def\){\right)} 
\newcommand\pt{p_{\sss\rm T}}
\newcommand{\ii}{\mathord{\mathrm{i}}}
\newcommand{\vbfnlo}{\textsc{VBFNLO}}
\newcommand{\bases}{\textsc{BASES}}
\newcommand{\mulbos}{\textsc{MulBos}}
\newcommand{\formcalc}{\textsc{FormCalc}}
\newcommand{\feynarts}{\textsc{FeynArts}}
\newcommand{\madgraph}{\textsc{MadGraph}}
\newcommand{\sherpa}{\textsc{Sherpa}}
\newcommand{\mocanlo}{\textsc{MoCaNLO}}
\newcommand{\loopints}{\textsc{LoopInts}}
\newcommand{\looptools}{\textsc{LoopTools}}
\newcommand{\lhapdf}{\textsc{LHAPDF}}
\newcommand{\msbar}{$\overline{\text{MS}}$}
\newcommand{\sect}[1]{Section~\ref{#1}}
\newcommand{\tab}[1]{Table~\ref{#1}}
\newcommand{\eq}[1]{Eq.~(\ref{#1})}
\newcommand{\bib}[1]{Ref.~\cite{#1}}
\newcommand{\refs}[1]{Refs.~\cite{#1}}
\newcommand{\bibs}[1]{\cite{#1}}
\newcommand{\fig}[1]{Fig.~\ref{#1}}
\newcommand{\Openmpi}{\mbox{\textsc{Openmpi}}}

%\textwidth=16cm \textheight=22.5 cm \topmargin 0cm \oddsidemargin-0.2cm

\geometry{verbose,a4paper,tmargin=2.5cm,bmargin=2.5cm,lmargin=2.5cm,rmargin=2.5cm}

% rotated table headers
\newcommand{\rot}[1]{\rlap{\rotatebox{45}{#1}}\hspace*{1em}}
\newcommand{\rotl}[1]{{\rotatebox{45}{#1}}}
\newcommand{\mcrot}[4]{\multicolumn{#1}{#2}{\rlap{\rotatebox{#3}{#4}~}}} 
% test letter in string
\makeatletter
\def\instring#1#2{TT\fi\begingroup
  \edef\x{\endgroup\noexpand\in@{#1}{#2}}\x\ifin@}
\makeatother
\newcommand{\makeflag}[3]{%
\if\instring{#1}{#3}{$\checkmark$}\else\if\instring{#2}{#3}{$\bigcirc$}\else{$-$}\fi\fi%
}
\newcommand{\bsmoptionsBLHA}[1]{%
\makeflag{B}{b}{#1} % BLHA
& 
\makeflag{L}{l}{#1} % semi-leptonic decay
&
\makeflag{F}{f}{#1} % VBF process
&
\makeflag{V}{v}{#1} % anomalous Gauge couplings
&
\makeflag{H}{h}{#1} % anomalous Higgs couplings
&
\makeflag{T}{t}{#1} % Two-Higgs model
&
\makeflag{K}{k}{#1} % Kaluza-Klein model
&
\makeflag{S}{s}{#1} % Spin-2 model
&
\makeflag{M}{m}{#1} % MSSM
}


\newcommand{\bsmgfoptionsBLHA}[1]{%
\makeflag{B}{b}{#1} % BLHA
&
\makeflag{G}{g}{#1} % gluon-fusion process
&
\makeflag{L}{l}{#1} % semi-leptonic decay
&
\makeflag{H}{h}{#1} % anomalous Higgs couplings
&
\makeflag{T}{t}{#1} % general 2HDM
&
\makeflag{M}{m}{#1} % MSSM
}
\newcommand{\bsmoptions}[1]{%
\makeflag{L}{l}{#1} % semi-leptonic decay
&
\makeflag{F}{f}{#1} % VBF process
&
\makeflag{V}{v}{#1} % anomalous Gauge couplings
&
\makeflag{H}{h}{#1} % anomalous Higgs couplings
&
\makeflag{T}{t}{#1} % Two-Higgs model
&
\makeflag{K}{k}{#1} % Kaluza-Klein model
&
\makeflag{S}{s}{#1} % Spin-2 model
&
\makeflag{M}{m}{#1} % MSSM
}
\newcommand{\bsmgfoptions}[1]{%
\makeflag{G}{g}{#1} % gluon-fusion process
&
\makeflag{L}{l}{#1} % semi-leptonic decay
&
\makeflag{H}{h}{#1} % anomalous Higgs couplings
&
\makeflag{T}{t}{#1} % general 2HDM
&
\makeflag{M}{m}{#1} % MSSM
}


\begin{document}
\begin{titlepage}
{\begin{flushright}{
 \begin{minipage}{8.5cm}
  %Phenikaa/PIAS/2025/XX
\end{minipage}}\end{flushright}}
\vspace{1cm}
\begin{center} {\Large \bf \textsc{MulBos}: A Monte Carlo program
    } 
\end{center}
\begin{center}
 {\Large \bf  for polarized multiboson production processes --}
\end{center}
\begin{center}
 {\Large \bf  Manual for Version 1.0}
\end{center}
\vspace{0.6cm}
\begin{center}

{\renewcommand{\baselinestretch}{4}
Duc Ninh Le$^{1}$,
Thi Nhung Dao$^{1}$,
Van Cuong Le$^{1}$
%M.~L\"oschner$^{7}$
}
\end{center}
\vspace{0.4cm}
\begin{center}
  $^{1}$ Phenikaa Institute for Advanced Study, Phenikaa University, Hanoi 12116, Vietnam \\ \noindent
  %$^{9}$ Particle  Physics,  Faculty  of  Physics,  University  of  Vienna,  Austria\\ \noindent  
\noindent
\end{center}
\vspace{0.6cm}

\newpage

\begin{abstract}

\noindent
  \textsc{MulBos} (for MultiBoson production) is a parton-level program to calculate polarized cross sections 
  for $W^+W^-$, $W^\pm Z$, and $ZZ$ production processes with fully leptonic decays at the LHC. 
  The calculations are performed at next-to-leading order (NLO) QCD+EW accuracy. 
  In this manual, we explain how to install and run the program. Important practical details are also provided. 
  Extension to triboson production processes will be added in the future. 
\end{abstract}
%\vspace{2 cm}
\today
\end{titlepage}

\newpage

\tableofcontents

%--------------------------------------------------------------------------------
%================= Section ======================================================
%--------------------------------------------------------------------------------
\newpage

%\section{Introduction}
%TBF.
%%%
\section{Installation and run}
The source code of \mulbos\ can be downloaded at

\begin{center}
\url{https://github.com/mulbos/mulbos}
\end{center}

%===================================================================

\subsection{Prerequisites}

The basic installation requires GNU {\tt make}, the {\sc Fortran} compiler Intel {\tt ifx}. 
A {\tt C} compiler is required for \loopints , which has been 
tested with {\tt gcc}.
The library \lhapdf \footnote{\tt
\url{https://www.lhapdf.org/}}~\cite{Buckley:2014ana} is needed
for parton distribution functions (PDF). 

For easy compilation and run of the program, the simple python scripts {\tt setup.py} 
and {\tt run\_VV.py} (with {\tt VV=WZ, WpWm, ZZ}, located in {\tt run\_templates}) are provided. 
To run these scripts, {\tt python3} is needed. 
Please note that \lhapdf\ also requires {\tt python3} for installation. 
%===================================================================

\subsection{Compilation and run}
\label{sec:compile}
Upon downloading \mulbos, one goes to the home directory of the program. 
From here, the compilation can be done with one command: 
\begin{center}
{\tt python3 setup.py}
\end{center}
This command installs all necessary libraries, except for the \lhapdf. 
\lhapdf\ (version $\ge$ 6.2.3 is recommended) must be installed separately 
and the path to the installation folder must be provided in the {\tt run\_VV.py}.

The current version of \mulbos\ is capable of calculating all polarized cross sections 
for the 4 diboson processes $ZZ$, $W^+W^-$, $W^+Z$, and $W^-Z$ at NLO QCD+EW at proton-proton colliders. 
Any kinematic distribution can be generated. Event generation is not yet available.

Instructions to calculate polarized cross sections at NLO are given in \sect{sec:cal_pol_xs}. 
Here we just explain how to perform a run. 

To perform a calculation for a specific process, say $ZZ$, one first creates a personal working directory 
named e.g. {\tt my\_working}. This directory can be placed anywhere. A good place is the \mulbos\ home directory.
One then copy the folder {\tt run\_templates/ZZ} to this working directory.
Go to the folder {\tt my\_working/ZZ} and edit the file {\tt run\_ZZ.py} to
provide the path to \mulbos\ home directory (search {\tt home\_mulbos}) 
and the path to LHAPDF library (search {\tt LHAPDF}). Then execute the following commands:
\begin{center}
{\tt python3 run\_ZZ.py}
\end{center}
You can change the run parameters in {\tt def run()}, the input parameters in\\ 
{\tt cards/input.dat.in}, 
and the cut parameters in {\tt cards/cuts.dat.in}. 
These input and cut parameters have been chosen to be the same as \bib{Carrivale:2025mjy} for the sake of comparison.

Upon running the {\tt run\_ZZ.py} script, a new folder named {\tt results\_*} (which is the default name, you can change it via 
the variable {\tt output\_dir} in {\tt run\_ZZ.py}) is created if 
it does not yet exist. The results are then 
stored in this directory. Folders named {\tt run\_IDp\_IDc\_polXY\_*} are created \footnote{One can easily change this name in {\tt run\_ZZ.py}.}, 
where {\tt IDp} and {\tt IDc} are the process and contribution identification numbers (ID), respectively.  
The list of these identification numbers are provided in Tables \ref{tab:IDs_proc_cont_WpZ}, 
\ref{tab:IDs_proc_cont_WmZ}, \ref{tab:IDs_proc_cont_WW_U}, \ref{tab:IDs_proc_cont_WW_D}, 
\ref{tab:IDs_proc_cont_WW_gg_aa}, \ref{tab:IDs_proc_cont_WW_B}, 
\ref{tab:IDs_proc_cont_ZZ_U}, \ref{tab:IDs_proc_cont_ZZ_D}, 
\ref{tab:IDs_proc_cont_ZZ_gg_aa}, \ref{tab:IDs_proc_cont_ZZ_B}. 
The last index {\tt polXY} gives the polarization state, see \tab{tab:IDs_polars} for the list of 
polarization IDs ({\tt XY=IDh}). 

In the {\tt run\_IDp\_IDc\_polXY\_*} folder, the following files are created:
\begin{enumerate}
{\tt
\item makefile
\item main.F
\item input.dat
\item cuts.dat
\item mulbos
\item xs.txt
\item dist/
}
\end{enumerate}
This {\tt run\_IDp\_IDc\_polXY\_*} directory is a complete program in itself. 
One can copy this folder to another place and compile it to create the executable {\tt mulbos}, 
and then submit job from there. 
For an easy use, this compilation and job submission steps have been done in {\tt run\_ZZ.py}. 
After the job has finished, the integrated cross section is provided in the file {\tt xs.txt}, where
some basic information about the input parameters and integration parameters are also given. 
Kinematic distributions are provided in the {\tt dist} folder. 
The current version of the program cannot generate events.

To get the full NLO results, one has to sum over all relevant processes (e.g. $\bar{U}U$, $\bar{D}D$, $gg$, $\gamma\gamma$, $b\bar{b}$) 
and contributions (LO, virtual corrections, real emission corrections, ...) 
for each distribution (or integrated cross section). See \sect{sec:cal_pol_xs} for specific instructions. 
%===================================================================
\subsection{Bug reports}
If you encounter any problems, please inform us via {\tt leducninh@gmail.com} 
or open an issue in github: \url{https://github.com/mulbos/mulbos/issues}.
%%%
\section{Source and directory layout}
The \mulbos\ source tree contains the following subdirectories:
\begin{itemize}
\item {\tt driver/:} Containing the main program {\tt main.F} and a makefile template {\tt makefile.in} 
which is used to create the final makefile (located in {\tt run\_IDp\_IDc\_polXY} as above discussed).
\item {\tt tools/:}
\begin{itemize}
\item {\tt utilities/:} Containing routines for calculating model parameters, run 
parameters, 
generating momenta of external-state particles, 
on-shell mappings, kinematic cuts, PDFs, filling histograms of kinematic distributions, 
the integrands of various contributions (LO, virtual corrections, real-emission corrections, ...). 
The Catani-Seymour-Dittmaier dipole subtraction method is here implemented. 
Renormalization constants for virtual corrections are here included. 
Routines to read in input files and print out important information are placed here as well.
\item {\tt FormCalc\_Tools/:} These interface files are needed for the calculation of the \formcalc\ matrix 
elements in the {\tt processes} directory. They have been adapted from \formcalc-6.0. 
In more details, the purpose of these routines are to calculate spinors, polarization vectors for each external-state 
particle, and some basic functions of these quantities. 
\item {\tt integrators/:} Containing the Monte-Carlo (MC) integrator {\tt BASES-5.1} \cite{Kawabata:1985yt,Kawabata:1995th}, 
which is needed for the phase-space integration and kinematic distribution generation. 
\item {\tt loops/:} Containing the library \loopints, for the calculation of 
one-loop integrals needed for the virtual corrections or loop-induced processes.
\end{itemize}
%%%
\item {\tt processes/:} Routines to calculate helicity amplitudes for the processes
provided. The production and decay helicity 
amplitudes are generated by \formcalc\ + \feynarts. 
Additional work is needed to combine these amplitudes according to the PA 
to obtained multi-gauge boson polarized amplitudes. 
\formcalc-6.0 and \feynarts-3.4 have been used to generate all amplitudes for the diboson processes. 
Importantly, the list of processes, contributions and their IDs are provided in the {\tt processes\_list.py} file 
located in this folder.
\item {\tt run\_templates/:} Run templates for all processes.
\item {\tt manual/:} The Latex source files of this manual.
\end{itemize} 

%===================================================================
\section{License}
\mulbos\ is distributed under the GNU General Public License~(GPL)
version~3. This ensures that the source code will be available to users,
grants them the freedom to use and modify the program and sets
out the conditions under which it can be redistributed. More details 
are provided in the {\tt COPYING} file shipped together with the source code.
%%%
\section{Credits}
\mulbos\ contains codes taken and adapted from other open-source programs.
The matrix elements are based on automatically
generated code from {\sc FeynArts+FormCalc}~\cite{Kublbeck:1990xc, 
Hahn:1998yk, Hahn:2000kx, Hahn:2006qw}. 
Some main interface routines for process and input-parameter initialization are adapted 
from {\sc FormCalc-6.0} codes. 
The phase space generation codes, which generate off-shell momenta from random numbers, 
are based on codes taken from \vbfnlo\ \cite{Arnold:2008rz,Baglio:2014uba,Baglio:2024gyp}. 
The routines which read in input parameter values from the card files are 
also taken from \vbfnlo.

A very slightly modified version of the \bases-5.1 \cite{Kawabata:1985yt,Kawabata:1995th} library is shipped 
together with the code. This is needed for Monte-Carlo integration and the generation of kinematic distributions. 

One-loop integrals are calculated by the library \loopints. Here, some codes are 
taken and adapted from \looptools\ \cite{Hahn:1998yk}. Specifically, they include the cache system (see {\tt cache.c}), 
Passarino-Veltman reduction for 3-point and 4-point tensor integrals. The routines to solve the system 
of linear equations needed by this reduction method are also taken from \looptools, see {\tt solve-LU.F}. 
The core parts of \loopints, including 2-point integrals, 3-point and 4-point scalar integrals have been implemented by us from scratch. 
These codes have been cross checked and used in many calculations \cite{Nhung:2009pm,Boudjema:2009pw,Nhung:2013jta,Campanario:2013qba,Baglio:2018rcu,Dao:2024ffg}.
%%%
\section{Code flow}
\label{sec:code_flow}
After executing {\tt python run\_VV.py}, a new run folder named {\tt run\_*} is 
created in the {\tt results\_*} directory. 

Go to this run folder, one finds the main program {\tt main.F}, which reads in the input files 
({\tt input.dat} and {\tt cuts.dat} located in the same directory), 
initializes important run parameters (e.g. colliding energy, model parameters) 
and performs the phase space integration. There is a {\tt makefile} there, which links the 
the main program with the needed libraries.

The phase space integration, which is done via 
the subroutine {\tt IntegratedCS} located in {\tt tools/utilities/xsection.F}. 
Here, the number of integration variables is automatically calculated and 
the subroutine {\tt BASES2} is called. 
This is an interface to the program {\tt BASES} \cite{Kawabata:1995th}, 
which is provided in {\tt tools/integrators/basesv5.1}. 
{\tt BASES} reads in the integrand (named {\tt FUNC}) and then 
performs the MC integration. 
It returns the integrated cross section (the mean value and the standard deviation) 
and a set of kinematic distributions in the folder {\tt dist}. 

The value of the integrand {\tt FUNC} is calculated from the list of random variables $x(i)$, 
which are generated by {\tt BASES}. 
The main subroutine here is {\tt DifferentialCS} (located in {\tt tools/utilities/Integrand.F}) where all 
the external momenta, PDFs, helicity amplitudes are calculated. 
Kinematic cuts and histogram filling are done via this subroutine as well.
%%%
\section{Input files and parameters}
\label{sec:input} 
\mulbos\ reads in input parameters through the following input files:
\begin{itemize}
\item {\tt run\_VV.py}: Some run parameters.
\item {\tt input.dat.in}: Run and model parameters.
\item {\tt cuts.dat.in}: Kinematic cut parameters.
\end{itemize}
%%%
\section{Checks}
Extensive checks for all components have been done by consistency checks (e.g. UV and IR finiteness) 
as well as comparing with external programs (\vbfnlo, \sherpa \cite{Hoppe:2023uux}, \mocanlo \cite{Denner:2026phn}). 
%------------------------------------------------------------------------
%================= Section ==============================================
%------------------------------------------------------------------------
%\newpage
\section{Calculating polarized cross sections and distributions}
\label{sec:cal_pol_xs}
%%%
Here we give instructions to calculate polarized cross sections and distributions at NLO QCD+EW.
The list of Polarization ID numbers, denoted as IDh (h for helicity), are provided in \tab{tab:IDs_polars}. 
The list of Process ID (IDp) and Contribution ID (IDc) numbers are provided in 
\tab{tab:IDs_proc_cont_WpZ} for $W^+Z$; 
\tab{tab:IDs_proc_cont_WmZ} for $W^-Z$;
Tables~\ref{tab:IDs_proc_cont_WW_U},\ref{tab:IDs_proc_cont_WW_D},\ref{tab:IDs_proc_cont_WW_gg_aa}, 
\ref{tab:IDs_proc_cont_WW_B} for $W^+W^-$;
Tables~\ref{tab:IDs_proc_cont_ZZ_U},\ref{tab:IDs_proc_cont_ZZ_D},\ref{tab:IDs_proc_cont_ZZ_gg_aa}, 
\ref{tab:IDs_proc_cont_ZZ_B} for $ZZ$.
%%%
\begin{table}[th!]
% \renewcommand{\arraystretch}{1.3}
\begin{center}
%\footnotesize
%\setlength\tabcolsep{0.03cm}
%\fontsize{7.0}{7.0}
\begin{tabular}{|cc|cc|cc|cc|cc|cc|}\hline
IDh  & Name & IDh & Name & IDh & Name & IDh & Name & IDh & Name & IDh & Name \\
\hline
44 & UU & 22 & LL & 25 & LT & 41 & UM & 14 & MU & 45 & UT \\
52 & TL & 55 & TT & 24 & LU & 42 & UL & 24 & LU & 54 & TU \\
42 & UL & 54 & TU & 45 & UT & 43 & UP & 34 & PU & $\lambda_1\lambda_2$ &  \\
\hline
\end{tabular}
%%%
\caption{\small Polarization ID (IDh) numbers and names for the $V_1V_2$ processes. Longitudinal polarization 
is denoted by the letter L and number $2$, transverse polarization T and $5$, unpolarized case U and $4$, 
transverse left polarization M and $1$, transverse right polarization P and $3$. 
The first character is for the first gauge boson $V_1$, the second character for $V_2$. 
The last IDh denoted as $\lambda_1\lambda_2$ with $\lambda_i = 1,2,3$ is for other possibilities such as $12$, $13$, etc.}
\label{tab:IDs_polars}
\end{center}
\end{table}
%%%
\subsection{$WZ$}
\label{sec:cal_xs_WZ}
We first take the $W^+Z$ process (IDp = 1) in \tab{tab:IDs_proc_cont_WpZ}.
In order to obtain results at NLO QCD+EW one has to sum over all contributions with IDc from 1 to 17. 
Note that the special case of IDc = 0 is reserved for the unpolarized off-shell LO contribution. 
This is not needed for calculating polarized cross sections. It is however useful for crosschecking with other programs.   

Using the provided python script {\tt run\_WZ.py} in {\tt run\_templates}, one calculates all those 17 contributions for all polarizations as follows:
\begin{verbatim}
CALCULATION = "POLAR_NLOQCDEW"
IDp_range = [1]
IDh_range = [22,55,52,25,44]
\end{verbatim}
This script will create 17 output directories, 
and submit 17 jobs. After these jobs are finished, one then performs the sum over the results (either 
the integrated cross section or a given distribution) to obtain the desired NLO QCD+EW result.
For the $W^-Z$ process (see \tab{tab:IDs_proc_cont_WmZ}), one sets {\tt IDp\_range = [2]}. To submit jobs for both $W^+Z$ and $W^-Z$ processes, 
then {\tt IDp\_range = [1,2]}. 

To obtain the off-shell LO results, one sets {\tt CALCULATION = "off\_shell"}. Other options including {\tt "POLAR\_LO"}, {\tt "POLAR\_NLOQCD"} 
are available as given in {\tt run\_WZ.py}.
%%%
\subsection{$W^+W^-$}
\label{sec:cal_xs_WW}
The processes and contributions for the $W^+W^-$ case are provided in Tables \ref{tab:IDs_proc_cont_WW_U}, \ref{tab:IDs_proc_cont_WW_D}, 
\ref{tab:IDs_proc_cont_WW_gg_aa}, \ref{tab:IDs_proc_cont_WW_B}. 
They can all be calculated by invoking {\tt python3 run\_WpWm.py} with the following settings:
\begin{verbatim}
CALCULATION = "POLAR_NLOQCDEW"
# sub-dominant contributions for POLAR calculations:
gg_process = True
aa_process = True
bb_process = 2 # (choose 0 for NO bb, 1 for LO, 2 for NLO)
IDh_range = [22,55,52,25,44]
\end{verbatim}
After all jobs are finished, one then combines the results as follows. 
For contributions in Tables \ref{tab:IDs_proc_cont_WW_U}, \ref{tab:IDs_proc_cont_WW_D}, 
\ref{tab:IDs_proc_cont_WW_gg_aa}, we add them together as usual (the errors are combined in quadrature). 
For the bottom-induced contributions in \tab{tab:IDs_proc_cont_WW_B} we add the contributions with {\tt IDc} from 1 to 17. 
Combining them all together gives the $\text{YesTW}$ results in \bib{Dao:2024ffg}.
To obtain the $\text{NoTW}$ results in \bib{Dao:2024ffg} one then subtracts from the $\text{YesTW}$ results
the contributions with {\tt IDc} from 18 to 21 in \tab{tab:IDs_proc_cont_WW_B}.   
Note that the sum from 1 to 17 in \tab{tab:IDs_proc_cont_WW_B} includes the $WW$, $tW$ and their interferences, 
while the contributions from 18 to 21 are the on-shell $tW$.
%%%
\subsection{$ZZ$}
\label{sec:cal_xs_ZZ}
The processes and contributions for the $ZZ$ case are provided in Tables \ref{tab:IDs_proc_cont_ZZ_U}, \ref{tab:IDs_proc_cont_ZZ_D}, 
\ref{tab:IDs_proc_cont_ZZ_gg_aa}, \ref{tab:IDs_proc_cont_ZZ_B}. 
They can all be calculated by invoking {\tt python3 run\_ZZ.py} with the following settings:
\begin{verbatim}
CALCULATION = "POLAR_NLOQCDEW"
# sub-dominant contributions for POLAR calculations:
gg_process = True
bb_process = 2 # (choose 0 for NO bb, 1 for LO, 2 for NLO)
IDh_range = [22,55,52,25,44]
\end{verbatim}
The $\gamma\gamma$ process contributes only to the off-shell case and is included there. 
The NLO QCD+EW cross sections are then obtained by adding all contributions. 
%%%
\subsection{Practical advice}
\label{sec:advice}
%%% optimal statistics (2^15 to 2^25)
The statistics for each contribution should be adjusted according to the importance of the contribution. 
The importance level is judged based on the value of the integrated cross section and the impact on the final distributions.  
For the $VV$ at NLO QCD+EW, an optimal set of statistics for the up or down contributions from 1 to 17 reads
\begin{align}
\log_2 (N_\text{u,VV}^{1\to 17})=\{24,19,19,19,24,22,22,17,17,17,23,21,21,17,21,17,21\},
\end{align}
where the LO term receives the highest statistics due to its large cross section.
The \{5,6,7\} (for QCD) and \{11,12,13,15,17\} (for EW) contributions also receive significant statistics because 
of their impact on the distributions. There are significant fluctuations in the distributions of these contributions 
due to the cancellation nature of the dipole-subtraction method \cite{Catani:1996vz,Dittmaier:1999mb}. 
To smooth out these fluctuations, additionally, one needs to combine results of different random seeds, more below.

For the $gg$ and $\gamma\gamma$, a good choice is $N_\text{gg,VV} = 2^{17}$ for both $WW$ and $ZZ$. 

For the $b$-induced contributions of the $ZZ$ case one can choose $N_\text{b,ZZ}^{1\to 17}=N_\text{u,VV}^{1\to 17}/k$ 
with $k=2$.

For the $b$-induced contributions of the $WW$ case one can choose
\begin{align}
\log_2 (N_\text{b,WW}^{1\to 21})=\{23,18,18,18,23,23,23,16,16,16,22,22,22,16,20,16,20,23,23,21,21\},
\end{align}
where the QCD contributions \{6,7\} and \{18,19\} receive large statistics because they are both huge (containing 
a top resonance) and cancel each other in the $\text{NoTW}$ calculation. 
Similarly for the EW contributions \{12,13\} and \{20,21\}, with smaller statistics compared to their QCD counterparts 
due to smaller cross sections.

%%% combine seeds
To further reduce fluctuations in the kinematic distributions, a good strategy is combining results of different 
random seeds (one can start with 5 seeds, and increase it if needed). Two common methods for combination are 
the normal average (using equal weights) and the weighted mean (using weights proportional to $1/\Delta_i^2$ with 
$\Delta_i$ being the error of the $i$th seed). The weighted-mean method gives smaller error, but it can only 
be trusted when there are no biases in the calculation of $\Delta_i$. A good practice is to implement both methods and compare 
the results.

%%% scale uncertainties: mu_F and mu_R dependence
The standard way to calculate the scale uncertainties is to vary $(\mu_F, \mu_R)$ around the central values 
$(\mu_{F0}, \mu_{R0})$, up and down by a factor of $2$. A well-known method is the seven-point method. 
The seven scale points are $(\mu_F, \mu_R) = (i\mu_{F0}, j\mu_{R0})$ 
with $i=0.5,1,2$ and $j=0.5,1,2$, with the constraint of $i/j < 4$ and $j/i < 4$ to avoid wide-separated scales. 
In this calculation one can reduce significantly computing time by noting that only the QCD-correction 
contributions ({\tt IDc=2,3,4,5,6,7} for all $VV$ processes, for the case of $WW$ there are 
additionally {\tt IDc=18,19}) depend on both $\mu_F$ and $\mu_R$. The other contributions (LO and NLOEW corrections) 
depend only on $\mu_F$.  
%%%
\section{Guidelines for implementing new cuts}
\label{sec:cuts}
To implement a new kinematic cut, one first adds an entry to {\tt cards/cuts.dat.in} to give it a name. 
Then edit the subroutine {\tt Read\_Cuts()} in {\tt utilities/global\_cuts.F90} to read in this value. 
Then edit the subroutine {\tt DoCut\_DiBoson()} in {\tt utilities/Cut\_tools.F} to implement this cut. 
For the NLOEW contributions with a photon radiation, one must do charged lepton photon recombination before applying the cut.
Finally, one must run {\tt make clean} and {\tt make} in {\tt tools/utilities} to update the {\tt utilities} library.

A practical advice: if you are not sure about the momenta, 
you can use \\
{\tt vext\_jend(5,nlegs\_max)} to verify. 
{\tt vext\_jend} contains the information of the 4-momentum (first 4 entries)
and particle-ID (fifth entry) of all external particles. The particle-IDs are given here \cite{ParticleDataGroup:1998hll}, which is 
the standard particle numbering scheme for MC generators. 
Note that lepton-photon recombination has not been done on {\tt vext\_jend}.
%%%
\section{Guidelines for implementing new distributions}
\label{sec:dists}
To implement a new kinematic distribution, one first adds an entry to \\
{\tt utilities/Hist\_init\_VV.F} (with {\tt VV=WW,WZ,ZZ})
to initialize a histogram template. 
Notice there that each histogram has an ID whose value is given by the variable {\tt Ncout}. 
This ID will be used to fill the histogram in the corresponding {\tt utilities/Hist\_fill\_VV.F}. 
As in the case of cuts, one must do charged lepton photon recombination before calculating the bin-position ({\tt pos\_bin}) for 
filling the histogram \\
(command {\tt CALL XHFILL(ID,pos\_bin,weight)}).
%%%
%\section{Conclusions and Outlook}
%\section*{Acknowledgments}
%%%
\newpage
\appendix
%%% Process IDs %%%
\section{List of processes and contributions}
\label{sec:list_proc_cont}

\subsection{WZ}
\begin{table}[th!]
\begin{center}
\begin{tabular}{|c|c|c|c|c|}\hline
IDp & IDc & Name & Correction Type & Polarized \\
\hline
1 & 0 & {\tt UDbarWpZ4lfull-born} & LO & No\\
1 & 1 & {\tt UDbarWpZ-born} & LO & Yes\\
1 & 2 & {\tt UDbarWpZ-virtQCD} & NLOQCD & Yes\\
1 & 3 & {\tt UDbarWpZrad-collQCD} & NLOQCD & Yes\\
1 & 4 & {\tt UDbarWpZind-collQCD} & NLOQCD & Yes\\
1 & 5 & {\tt UDbarWpZglu-fin} & NLOQCD & Yes\\
1 & 6 & {\tt DbargluUbarWpZ-fin} & NLOQCD & Yes\\
1 & 7 & {\tt UgluDWpZ-fin} & NLOQCD & Yes\\
1 & 8 & {\tt UDbarWpZ-virtEW} & NLOEW & Yes\\
1 & 9 & {\tt UDbarWpZrad-collEW} & NLOEW & Yes\\
1 & 10 & {\tt UDbarWpZind-collEW} & NLOEW & Yes\\
1 & 11 & {\tt UDbarWpZgam-fin} & NLOEW & Yes\\
1 & 12 & {\tt DbargamUbarWpZ-fin} & NLOEW & Yes\\
1 & 13 & {\tt UgamDWpZ-fin} & NLOEW & Yes\\
1 & 14 & {\tt UDbarWpZ-virtEW-decayZ} & NLOEW & Yes\\
1 & 15 & {\tt UDbarWpZgam-fin-decayZ} & NLOEW & Yes\\
1 & 16 & {\tt UDbarWpZ-virtEW-decayW} & NLOEW & Yes\\
1 & 17 & {\tt UDbarWpZgam-fin-decayW} & NLOEW & Yes\\
\hline
\end{tabular}
\caption{\small Process ID (IDp) and Contribution ID (IDc) numbers for the $W^+Z$ process $pp \to e^+\nu_e\mu^+\mu^-$.}
\label{tab:IDs_proc_cont_WpZ}
\end{center}
\end{table}

\begin{table}[th!]
\begin{center}
\begin{tabular}{|c|c|c|c|c|}\hline
IDp & IDc & Name & Correction Type & Polarized \\
\hline
2 & 0 & {\tt UbarDWmZ4lfull-born} & LO & No\\
2 & 1 & {\tt UbarDWmZ-born} & LO & Yes\\
2 & 2 & {\tt UbarDWmZ-virtQCD} & NLOQCD & Yes\\
2 & 3 & {\tt UbarDWmZrad-collQCD} & NLOQCD & Yes\\
2 & 4 & {\tt UbarDWmZind-collQCD} & NLOQCD & Yes\\
2 & 5 & {\tt UbarDWmZglu-fin} & NLOQCD & Yes\\
2 & 6 & {\tt UbargluDbarWmZ-fin} & NLOQCD & Yes\\
2 & 7 & {\tt DgluUWmZ-fin} & NLOQCD & Yes\\
2 & 8 & {\tt UbarDWmZ-virtEW} & NLOEW & Yes\\
2 & 9 & {\tt UbarDWmZrad-collEW} & NLOEW & Yes\\
2 & 10 & {\tt UbarDWmZind-collEW} & NLOEW & Yes\\
2 & 11 & {\tt UbarDWmZgam-fin} & NLOEW & Yes\\
2 & 12 & {\tt UbargamDbarWmZ-fin} & NLOEW & Yes\\
2 & 13 & {\tt DgamUWmZ-fin} & NLOEW & Yes\\
2 & 14 & {\tt UbarDWmZ-virtEW-decayZ} & NLOEW & Yes\\
2 & 15 & {\tt UbarDWmZgam-fin-decayZ} & NLOEW & Yes\\
2 & 16 & {\tt UbarDWmZ-virtEW-decayW} & NLOEW & Yes\\
2 & 17 & {\tt UbarDWmZgam-fin-decayW} & NLOEW & Yes\\
\hline
\end{tabular}
\caption{\small ID numbers for the $W^-Z$ process $pp \to e^-\bar{\nu}_e\mu^+\mu^-$.}
\label{tab:IDs_proc_cont_WmZ}
\end{center}
\end{table}

\newpage
\subsection{WW}

% (All WW tables updated in the same way: Contribution ID → ID−1)

\begin{table}[th!]
\begin{center}
\begin{tabular}{|c|c|c|c|c|}\hline
IDp & IDc & Name & Correction Type & Polarized \\
\hline
3 & 0 & {\tt UbarUWW4lfull-born} & LO & No\\
3 & 1 & {\tt UbarUWW-born} & LO & Yes\\
3 & 2 & {\tt UbarUWW-virtQCD} & NLOQCD & Yes\\
3 & 3 & {\tt UbarUWWrad-collQCD} & NLOQCD & Yes\\
3 & 4 & {\tt UbarUWWind-collQCD} & NLOQCD & Yes\\
3 & 5 & {\tt UbarUWWglu-fin} & NLOQCD & Yes\\
3 & 6 & {\tt UbargluUbarWW-fin} & NLOQCD & Yes\\
3 & 7 & {\tt UgluUWW-fin} & NLOQCD & Yes\\
3 & 8 & {\tt UbarUWW-virtEW} & NLOEW & Yes\\
3 & 9 & {\tt UbarUWWrad-collEW} & NLOEW & Yes\\
3 & 10 & {\tt UbarUWWind-collEW} & NLOEW & Yes\\
3 & 11 & {\tt UbarUWWgam-fin} & NLOEW & Yes\\
3 & 12 & {\tt UbargamUbarWW-fin} & NLOEW & Yes\\
3 & 13 & {\tt UgamUWW-fin} & NLOEW & Yes\\
3 & 14 & {\tt UbarUWW-virtEW-decayWm} & NLOEW & Yes\\
3 & 15 & {\tt UbarUWWgam-fin-decayWm} & NLOEW & Yes\\
3 & 16 & {\tt UbarUWW-virtEW-decayWp} & NLOEW & Yes\\
3 & 17 & {\tt UbarUWWgam-fin-decayWp} & NLOEW & Yes\\
\hline
\end{tabular}
\caption{\small ID numbers for the up-type quark induced $W^+W^-$ processes $pp(U) \to e^+\nu_e\mu^-\bar{\nu}_\mu$. 
Each contribution includes $u$ and $c$ induced processes.}
\label{tab:IDs_proc_cont_WW_U}
\end{center}
\end{table}

% Remaining WW, ZZ tables follow the same rule (all already shifted by −1)
\begin{table}[th!]
\begin{center}
\begin{tabular}{|c|c|c|c|c|}\hline
IDp & IDc & Name & Correction Type & Polarized \\
\hline
4 & 0 & {\tt DbarDWW4lfull-born} & LO & No\\
4 & 1 & {\tt DbarDWW-born} & LO & Yes\\
4 & 2 & {\tt DbarDWW-virtQCD} & NLOQCD & Yes\\
4 & 3 & {\tt DbarDWWrad-collQCD} & NLOQCD & Yes\\
4 & 4 & {\tt DbarDWWind-collQCD} & NLOQCD & Yes\\
4 & 5 & {\tt DbarDWWglu-fin} & NLOQCD & Yes\\
4 & 6 & {\tt DbargluDbarWW-fin} & NLOQCD & Yes\\
4 & 7 & {\tt DgluDWW-fin} & NLOQCD & Yes\\
4 & 8 & {\tt DbarDWW-virtEW} & NLOEW & Yes\\
4 & 9 & {\tt DbarDWWrad-collEW} & NLOEW & Yes\\
4 & 10 & {\tt DbarDWWind-collEW} & NLOEW & Yes\\
4 & 11 & {\tt DbarDWWgam-fin} & NLOEW & Yes\\
4 & 12 & {\tt DbargamDbarWW-fin} & NLOEW & Yes\\
4 & 13 & {\tt DgamDWW-fin} & NLOEW & Yes\\
4 & 14 & {\tt DbarDWW-virtEW-decayWm} & NLOEW & Yes\\
4 & 15 & {\tt DbarDWWgam-fin-decayWm} & NLOEW & Yes\\
4 & 16 & {\tt DbarDWW-virtEW-decayWp} & NLOEW & Yes\\
4 & 17 & {\tt DbarDWWgam-fin-decayWp} & NLOEW & Yes\\
\hline
\end{tabular}
\caption{\small ID numbers for the down-type quark induced $W^+W^-$ processes $pp(D) \to e^+\nu_e\mu^-\bar{\nu}_\mu$. 
Each contribution includes $d$ and $s$ induced processes.}
\label{tab:IDs_proc_cont_WW_D}
\end{center}
\end{table}

%%%%%% GGGGG %%%%%
\begin{table}[th!]
\begin{center}
\begin{tabular}{|c|c|c|c|c|}\hline
IDp & IDc & Name & Correction Type & Polarized \\
\hline
5 & 1 & {\tt glugluWW-loop} & NNLO & Yes\\
7 & 1 & {\tt gamgamWW-born} & LO & Yes\\
\hline
\end{tabular}
\caption{\small ID numbers for the gluon-gluon loop induced and photon-photon induced $W^+W^-$ processes contributing to $pp \to e^+\nu_e\mu^-\bar{\nu}_\mu$.}
\label{tab:IDs_proc_cont_WW_gg_aa}
\end{center}
\end{table}

%%%%%% MMMMM %%%%
\begin{table}[th!]
\begin{center}
\begin{tabular}{|c|c|c|c|c|}\hline
IDp & IDc & Name & Correction Type & Polarized \\
\hline
6 & 0 & {\tt BbarBWW4lfull-born} & LO & No\\
6 & 1 & {\tt BbarBWW-born} & LO & Yes\\
6 & 2 & {\tt BbarBWW-virtQCD} & NLOQCD & Yes\\
6 & 3 & {\tt BbarBWWrad-collQCD} & NLOQCD & Yes\\
6 & 4 & {\tt BbarBWWind-collQCD} & NLOQCD & Yes\\
6 & 5 & {\tt BbarBWWglu-fin} & NLOQCD & Yes\\
6 & 6 & {\tt BbargluBbarWW-fin} & NLOQCD & Yes\\
6 & 7 & {\tt BgluBWW-fin} & NLOQCD & Yes\\
6 & 8 & {\tt BbarBWW-virtEW} & NLOEW & Yes\\
6 & 9 & {\tt BbarBWWrad-collEW} & NLOEW & Yes\\
6 & 10 & {\tt BbarBWWind-collEW} & NLOEW & Yes\\
6 & 11 & {\tt BbarBWWgam-fin} & NLOEW & Yes\\
6 & 12 & {\tt BbargamBbarWW-fin} & NLOEW & Yes\\
6 & 13 & {\tt BgamBWW-fin} & NLOEW & Yes\\
6 & 14 & {\tt BbarBWW-virtEW-decayWm} & NLOEW & Yes\\
6 & 15 & {\tt BbarBWWgam-fin-decayWm} & NLOEW & Yes\\
6 & 16 & {\tt BbarBWW-virtEW-decayWp} & NLOEW & Yes\\
6 & 17 & {\tt BbarBWWgam-fin-decayWp} & NLOEW & Yes\\
6 & 18 & {\tt BbargluBbarWW-fin-TW} & NLOQCD & Yes\\
6 & 19 & {\tt BgluBWW-fin-TW} & NLOQCD & Yes\\
6 & 20 & {\tt BbargamBbarWW-fin-TW} & NLOEW & Yes\\
6 & 21 & {\tt BgamBWW-fin-TW} & NLOEW & Yes\\
\hline
\end{tabular}
\caption{\small ID numbers for the bottom-quark induced $W^+W^-$ processes $pp(B) \to e^+\nu_e\mu^-\bar{\nu}_\mu$.}
\label{tab:IDs_proc_cont_WW_B}
\end{center}
\end{table}


%%%%%% NNNNN %%%%
\newpage
\subsection{ZZ}

\begin{table}[th!]
\begin{center}
\begin{tabular}{|c|c|c|c|c|}\hline
IDp & IDc & Name & Correction Type & Polarized \\
\hline
8 & 0 & {\tt UbarUZZ4lfull-born} & LO & No\\
8 & 1 & {\tt UbarUZZ-born} & LO & Yes\\
8 & 2 & {\tt UbarUZZ-virtQCD} & NLOQCD & Yes\\
8 & 3 & {\tt UbarUZZrad-collQCD} & NLOQCD & Yes\\
8 & 4 & {\tt UbarUZZind-collQCD} & NLOQCD & Yes\\
8 & 5 & {\tt UbarUZZglu-fin} & NLOQCD & Yes\\
8 & 6 & {\tt UbargluUbarZZ-fin} & NLOQCD & Yes\\
8 & 7 & {\tt UgluUZZ-fin} & NLOQCD & Yes\\
8 & 8 & {\tt UbarUZZ-virtEW} & NLOEW & Yes\\
8 & 9 & {\tt UbarUZZrad-collEW} & NLOEW & Yes\\
8 & 10 & {\tt UbarUZZind-collEW} & NLOEW & Yes\\
8 & 11 & {\tt UbarUZZgam-fin} & NLOEW & Yes\\
8 & 12 & {\tt UbargamUbarZZ-fin} & NLOEW & Yes\\
8 & 13 & {\tt UgamUZZ-fin} & NLOEW & Yes\\
8 & 14 & {\tt UbarUZZ-virtEW-decayZ2} & NLOEW & Yes\\
8 & 15 & {\tt UbarUZZgam-fin-decayZ2} & NLOEW & Yes\\
8 & 16 & {\tt UbarUZZ-virtEW-decayZ1} & NLOEW & Yes\\
8 & 17 & {\tt UbarUZZgam-fin-decayZ1} & NLOEW & Yes\\
\hline
\end{tabular}
\caption{\small ID numbers for the up-type quark induced $ZZ$ processes $pp(U) \to e^+e^-\mu^+\mu^-$. 
Each contribution includes $u$ and $c$ induced processes.}
\label{tab:IDs_proc_cont_ZZ_U}
\end{center}
\end{table}

\begin{table}[th!]
\begin{center}
\begin{tabular}{|c|c|c|c|c|}\hline
IDp & IDc & Name & Correction Type & Polarized \\
\hline
9 & 0 & {\tt DbarDZZ4lfull-born} & LO & No\\
9 & 1 & {\tt DbarDZZ-born} & LO & Yes\\
9 & 2 & {\tt DbarDZZ-virtQCD} & NLOQCD & Yes\\
9 & 3 & {\tt DbarDZZrad-collQCD} & NLOQCD & Yes\\
9 & 4 & {\tt DbarDZZind-collQCD} & NLOQCD & Yes\\
9 & 5 & {\tt DbarDZZglu-fin} & NLOQCD & Yes\\
9 & 6 & {\tt DbargluDbarZZ-fin} & NLOQCD & Yes\\
9 & 7 & {\tt DgluDZZ-fin} & NLOQCD & Yes\\
9 & 8 & {\tt DbarDZZ-virtEW} & NLOEW & Yes\\
9 & 9 & {\tt DbarDZZrad-collEW} & NLOEW & Yes\\
9 & 10 & {\tt DbarDZZind-collEW} & NLOEW & Yes\\
9 & 11 & {\tt DbarDZZgam-fin} & NLOEW & Yes\\
9 & 12 & {\tt DbargamDbarZZ-fin} & NLOEW & Yes\\
9 & 13 & {\tt DgamDZZ-fin} & NLOEW & Yes\\
9 & 14 & {\tt DbarDZZ-virtEW-decayZ2} & NLOEW & Yes\\
9 & 15 & {\tt DbarDZZgam-fin-decayZ2} & NLOEW & Yes\\
9 & 16 & {\tt DbarDZZ-virtEW-decayZ1} & NLOEW & Yes\\
9 & 17 & {\tt DbarDZZgam-fin-decayZ1} & NLOEW & Yes\\
\hline
\end{tabular}
\caption{\small ID numbers for the down-type quark induced $ZZ$ processes $pp(D) \to e^+e^-\mu^+\mu^-$. 
Each contribution includes $d$ and $s$ induced processes.}
\label{tab:IDs_proc_cont_ZZ_D}
\end{center}
\end{table}

\begin{table}[th!]
\begin{center}
\begin{tabular}{|c|c|c|c|c|}\hline
IDp & IDc & Name & Correction Type & Polarized \\
\hline
10 & 1 & {\tt glugluZZ-loop} & NNLO & Yes\\
12 & 0 & {\tt gamgamZZ4lfull-born} & LO & No\\
\hline
\end{tabular}
\caption{\small ID numbers for the gluon-gluon loop induced and photon-photon induced $ZZ$ processes contributing to $pp \to e^+e^-\mu^+\mu^-$.}
\label{tab:IDs_proc_cont_ZZ_gg_aa}
\end{center}
\end{table}

\begin{table}[th!]
\begin{center}
\begin{tabular}{|c|c|c|c|c|}\hline
IDp & IDc & Name & Correction Type & Polarized \\
\hline
11 & 0 & {\tt BbarBZZ4lfull-born} & LO & No\\
11 & 1 & {\tt BbarBZZ-born} & LO & Yes\\
11 & 2 & {\tt BbarBZZ-virtQCD} & NLOQCD & Yes\\
11 & 3 & {\tt BbarBZZrad-collQCD} & NLOQCD & Yes\\
11 & 4 & {\tt BbarBZZind-collQCD} & NLOQCD & Yes\\
11 & 5 & {\tt BbarBZZglu-fin} & NLOQCD & Yes\\
11 & 6 & {\tt BbargluBbarZZ-fin} & NLOQCD & Yes\\
11 & 7 & {\tt BgluBZZ-fin} & NLOQCD & Yes\\
11 & 8 & {\tt BbarBZZ-virtEW} & NLOEW & Yes\\
11 & 9 & {\tt BbarBZZrad-collEW} & NLOEW & Yes\\
11 & 10 & {\tt BbarBZZind-collEW} & NLOEW & Yes\\
11 & 11 & {\tt BbarBZZgam-fin} & NLOEW & Yes\\
11 & 12 & {\tt BbargamBbarZZ-fin} & NLOEW & Yes\\
11 & 13 & {\tt BgamBZZ-fin} & NLOEW & Yes\\
11 & 14 & {\tt BbarBZZ-virtEW-decayZ2} & NLOEW & Yes\\
11 & 15 & {\tt BbarBZZgam-fin-decayZ2} & NLOEW & Yes\\
11 & 16 & {\tt BbarBZZ-virtEW-decayZ1} & NLOEW & Yes\\
11 & 17 & {\tt BbarBZZgam-fin-decayZ1} & NLOEW & Yes\\
\hline
\end{tabular}
\caption{\small ID numbers for the bottom-quark induced $ZZ$ processes $pp(B) \to e^+e^-\mu^+\mu^-$.}
\label{tab:IDs_proc_cont_ZZ_B}
\end{center}
\end{table}
%%%

\newpage
\section{Details of contributions}
\label{sec:cont_details}
%%%%% Details of contributions %%%%%
\subsection{\tt XX...4lfull-born}
\label{sec:full-born}
This is the full off-shell calculation for the $XX \to 4\ \text{leptons}$ process at LO. 
PDFs are included. Calculation of unpolarized cross sections can be performed. 
Polarized cross sections are unavailable.   
\subsection{\tt XX...V1V2-born}
\label{sec:DPA-born}
This is the DPA calculation for the $XX \to V_1 V_2 \to 4\ \text{leptons}$ process at LO. 
PDFs are included. Calculation of polarized cross sections can be performed.   
\subsection{\tt XX...V1V2-virtQCD}
\label{sec:virtQCD}
This is the DPA calculation for the QCD virtual corrections for the $XX \to V_1 V_2$ process at NLOQCD. 
Fully leptonic decays are implicitly understood. 
PDFs are included. Calculation of polarized cross sections can be performed. 
The QCD virtual corrections here include virtual loop corrections and soft-gluon real corrections 
(called endpoint contribution in \bib{Dittmaier:1999mb}). 
This contribution can be calculated using mass regularization or dimensional regularization. 
The default option is mass regularization.
\subsection{\tt XX...V1V2rad-collQCD}
\label{sec:rad-collQCD}
This is the DPA calculation for the so-called collinear corrections of the real gluon emission process 
$XX \to V_1 V_2 g$ occurring at NLOQCD. PDFs are included. 
Calculation of polarized cross sections can be performed. 
This contribution can be calculated using mass regularization or dimensional regularization. 
The default option is mass regularization.
This contribution is called the PK-operator contribution in \cite{Catani:1996vz}. 
The implementation of mass regularization is based on \cite{Dittmaier:1999mb}. 
PDF counter terms are here included. 
\subsection{\tt XX...V1V2ind-collQCD}
\label{sec:ind-collQCD}
This contribution is similar to the one described in \sect{sec:rad-collQCD}, but for the 
(anti)quark-gluon induced process $qg \to V_1 V_2 q^\prime$. 
\subsection{\tt XX...V1V2glu-fin}
\label{sec:glu-rad-fin}
This is the DPA calculation for the remaining finite corrections of the real gluon emission process 
$XX \to V_1 V_2 g$ occurring at NLOQCD. PDFs are included. 
Calculation of polarized cross sections can be performed. 

Dipole subtraction terms \cite{Dittmaier:1999mb} are added here such that the integration can be performed in 4 dimensions using 
massless quarks. 
\subsection{\tt XbargluYbar...V1V2-fin}
\label{sec:Qbarglu-ind-fin}
This contribution is similar to the one described in \sect{sec:glu-rad-fin}, but for the 
antiquark-gluon induced process $\bar{q}g \to V_1 V_2 \bar{q}^\prime$. 
\subsection{\tt XgluY...V1V2-fin}
\label{sec:Qglu-ind-fin}
This contribution is similar to the one described in \sect{sec:glu-rad-fin}, but for the 
quark-gluon induced process $qg \to V_1 V_2 q^\prime$. 
\subsection{\tt XX...V1V2-virtEW}
\label{sec:virtEW}
This contribution is similar to the one described in \sect{sec:virtQCD}, but for the 
EW corrections to the production part (i.e. $XX \to V_1 V_2$). 
\subsection{\tt XX...V1V2rad-collEW}
\label{sec:rad-collEW}
This contribution is similar to the one described in \sect{sec:rad-collQCD}, but for the 
EW corrections. 
The corrections here come from the real photon emission process 
$XX \to V_1 V_2 \gamma$.
PDF counter terms are here included, using the DIS scheme or $\overline{\text{MS}}$ scheme provided in \cite{Dittmaier:2009cr}. 
\subsection{\tt XX...V1V2ind-collEW}
\label{sec:ind-collEW}
This contribution is similar to the one described in \sect{sec:rad-collEW}, but for the 
(anti)quark-photon induced process $q\gamma \to V_1 V_2 q^\prime$.
\subsection{\tt XX...V1V2gam-fin}
\label{sec:gam-rad-fin}
This contribution is similar to the one described in \sect{sec:glu-rad-fin}, but for the 
real photon emission process $XX \to V_1 V_2 \gamma$.
\subsection{\tt XbargamYbar...V1V2-fin}
\label{sec:Qbargam-ind-fin}
This contribution is similar to the one described in \sect{sec:gam-rad-fin}, but for the 
antiquark-photon induced process $\bar{q}\gamma \to V_1 V_2 \bar{q}^\prime$.
\subsection{\tt XgamY...V1V2-fin}
\label{sec:Qgam-ind-fin}
This contribution is similar to the one described in \sect{sec:Qbargam-ind-fin}, but for the 
quark-photon induced process $q\gamma \to V_1 V_2 q^\prime$.
\subsection{\tt XX...V1V2-virtEW-decayVi}
\label{sec:virtEW-decay}
This contribution is similar to the one described in \sect{sec:virtEW}, but for the 
virtual EW corrections to the decay $V_i\to 2\,\text{leptons}$. 
The process is $XX \to V_1 V_2 \to 4\ \text{leptons}$.
\subsection{\tt XX...V1V2gam-fin-decayVi}
\label{sec:gam-rad-fin-decay}
This contribution is similar to the one described in \sect{sec:gam-rad-fin}, but for the 
real photon emission off the decay $V_i\to 2\,\text{leptons}$. 
The process is $XX \to V_1 V_2 \to 4\,\text{leptons}\,+\,\gamma$.
\subsection{\tt BbargluBbarWW-fin-TW}
\label{sec:Bbarg-ind-fin-TW}
This is an on-shell $tW$ contribution to the $W^+W^-$ process, originated from the $\bar{b}g$ induced 
initial state. This contribution is needed to subtract the on-shell $tW$ contribution occurring in 
the $\bar{b}g \to W^+W^- \bar{b}$ process (the {\tt BbargluBbarWW-fin} contribution 
described in \sect{sec:Qbarglu-ind-fin}). 
PDFs are included. Calculation of polarized cross sections can be performed. 
The process is $\bar{b}g \to W^+ \bar{t} \to W^+W^-\bar{b}$.
Details of this contribution are provided in \cite{Dao:2024ffg}.
\subsection{\tt BgluBWW-fin-TW}
\label{sec:Bg-ind-fin-qcd-TW}
This contribution is similar to the one described in \sect{sec:Bbarg-ind-fin-TW}, but originated 
from the $bg$ induced process $bg \to W^- t \to W^+W^-b$.
Details are also provided in \cite{Dao:2024ffg}.
\subsection{\tt BbargamBbarWW-fin-TW}
\label{sec:Bbargam-ind-fin-TW}
This contribution is similar to the one described in \sect{sec:Bbarg-ind-fin-TW}, but originated 
from the $\bar{b}\gamma$ induced process $\bar{b}\gamma \to W^+ \bar{t} \to W^+W^-b$.
Details are provided in \cite{Dao:2024ffg}.
\subsection{\tt BgamBWW-fin-TW}
\label{sec:Bgam-ind-fin-TW}
This contribution is similar to the one described in \sect{sec:Bbargam-ind-fin-TW}, but originated 
from the $b\gamma$ induced process $b\gamma \to W^- t \to W^+W^-b$.
Details are provided in \cite{Dao:2024ffg}.




\newpage


%--------------------------------------------------------------------------------
%================= Section ======================================================
%--------------------------------------------------------------------------------

\newpage

\bibliographystyle{JHEP}
\bibliography{MulBos}
\end{document}

%%%%%%%%%%%%%%%%%%
% other things:
%%%%%%%%%%%%%%%%%%
%%%
\section{Regularization for IR divergences}
To calculate NLO corrections (either QCD or EW), one must use a regularization method to deal 
with IR divergences. For NLO QCD, one can use either mass regularization or dimensional regularization. 
For NLO EW, mass regularization has been fully implemented. 
Dimensional regularization has been largely implemented, but there are still some missing corrections (e.g. I-operator for the $V$ decays). 
The default choice is therefore mass regularization, which is set via {\tt lMassReg = .true.} in {\tt global\_vars.F90}. 

\section{Off-shell contribution}
Off-shell contribution can be made small by imposing kinematic cuts targeting 
the double-resonance region. However, there can be some phase-space regions 
where off-shell effects are significant (being around $10\%$). 
These effects can be neglected in the first approximation. 
For precise measurements, they must be taken into account though.   

Off-shell contribution is 
the difference between the full off-shell and the unpolarized DPA cross sections, 
which should be calculated at the same perturbative order. 
The unpolarized DPA results (denoted as $44$ in \tab{tab:IDs_polars}) 
are provided by \mulbos\ at NLOQCDEW for quark-induced processes and at LO for 
the gluon-gluon and photon-photon processes. Only a few LO off-shell processes 
(denoted by keyword {\tt full} in 
Tables~\ref{tab:IDs_proc_cont_WpZ},\ref{tab:IDs_proc_cont_WmZ},\ref{tab:IDs_proc_cont_WW_U},\ref{tab:IDs_proc_cont_WW_D}...) 
are provided in our program. The NLOQCD results for the full off-shell 
processes can be obtained using e.g. \vbfnlo\ or \madgraph, while NLOEW corrections via e.g. 
\madgraph.   

%\section{Collaborations}
%Ideas for collaborations are welcome !
